% !TeX program = xelatex
% ============================================================================
%  第 8 课　课堂单页（学生用，单独打印一张 A4）
%  编译：XeLaTeX，两遍即可
%  约定：顶点 1 左上、2 右上、3 右下、4 左下；
%        r 顺时针转 90°，f 沿对角线 1—3 翻折；
%        复合 xy 读作「先做 x，再做 y」；格子里的动作是「先横排的，再做竖排的」
% ============================================================================

\documentclass[11pt]{ctexart}
\usepackage[a4paper,left=1.6cm,right=1.6cm,top=1.4cm,bottom=1.4cm]{geometry}
\usepackage{amsmath,amssymb}
\usepackage{booktabs}
\usepackage{array}
\usepackage[dvipsnames]{xcolor}

\setCJKmainfont{SimSun}[BoldFont=SimHei,ItalicFont=KaiTi]
\IfFontExistsTF{TeX Gyre Pagella}{\setmainfont{TeX Gyre Pagella}}{\setmainfont{Latin Modern Roman}}

\definecolor{pagegreen}{RGB}{34,139,34}
\renewcommand{\arraystretch}{1.2}
\setlength{\parindent}{0pt}
\pagestyle{empty}

\newcommand{\biaoti}[2]{%
  {\color{pagegreen}\rule{\textwidth}{1.2pt}}\\[6pt]
  {\heiti\Large 第 #1 课\quad #2}\hfill{\small\songti 课堂单页}\\[6pt]
  {\color{pagegreen}\rule{\textwidth}{0.5pt}}\\[4pt]
  {\small 姓名\underline{\hspace{3.2cm}}\qquad 班级\underline{\hspace{2.4cm}}\qquad 座号\underline{\hspace{1.6cm}}}
  \vspace{0.4em}
}

\begin{document}

\biaoti{8}{正方形对称：\(D_4\) 的 \(8\) 个元素}

\section*{一、八个动作的名字}

\begin{center}
{\small
\begin{tabular}{@{}lcc@{}}
  \toprule
  名字 & 动作 & 两行式 \\
  \midrule
  \(e\) & 不动 & \(\begin{pmatrix}1&2&3&4\\1&2&3&4\end{pmatrix}\) \\[3pt]
  \(r\) & 顺时针转 \(90^\circ\) & \(\begin{pmatrix}1&2&3&4\\2&3&4&1\end{pmatrix}\) \\[3pt]
  \(r^2\) & 顺时针转 \(180^\circ\) & \(\begin{pmatrix}1&2&3&4\\3&4&1&2\end{pmatrix}\) \\[3pt]
  \(r^3\) & 顺时针转 \(270^\circ\) & \(\begin{pmatrix}1&2&3&4\\4&1&2&3\end{pmatrix}\) \\[3pt]
  \(f\) & 沿对角线 \(1\text{—}3\) 翻 & \(\begin{pmatrix}1&2&3&4\\1&4&3&2\end{pmatrix}\) \\[3pt]
  \(rf\) & 沿\underline{\hspace{1.4em}}翻 & \(\begin{pmatrix}1&2&3&4\\\hspace{0.9em}&\hspace{0.9em}&\hspace{0.9em}&\hspace{0.9em}\end{pmatrix}\) \\[3pt]
  \(r^2f\) & 沿对角线 \underline{\hspace{1.4em}} 翻 & \(\begin{pmatrix}1&2&3&4\\3&2&1&4\end{pmatrix}\) \\[3pt]
  \(r^3f\) & 沿\underline{\hspace{1.4em}}翻 & \(\begin{pmatrix}1&2&3&4\\\hspace{0.9em}&\hspace{0.9em}&\hspace{0.9em}&\hspace{0.9em}\end{pmatrix}\) \\
  \bottomrule
\end{tabular}}
\end{center}

\section*{二、乘法表：只填 \(r\) 行与 \(f\) 行}

\noindent 格子里的动作是「\textbf{先横排的，再做竖排的}」。

\begin{center}
{\small
\begin{tabular}{@{}l|cccccccc@{}}
  \toprule
   & \(e\) & \(r\) & \(r^2\) & \(r^3\) & \(f\) & \(rf\) & \(r^2f\) & \(r^3f\) \\
  \midrule
  \rule{0pt}{3.2ex}先 \(r\) & & & & & & & & \\
  \rule{0pt}{3.2ex}先 \(f\) & & & & & & & & \\
  \rule{0pt}{3.2ex}先 \(r^2\) & & & & & & & & \\
  \bottomrule
\end{tabular}}
\end{center}

\noindent 拿纸片真的做两遍，或者用两行式一格一格推。

\section*{三、读出一个关系}

\noindent 看 \(f\) 行的第二个格子，它等于 \underline{\hspace{3em}}，所以可以写成一个等式：

\vspace{0.4em}
\noindent\hfill \(f\,r=\) \underline{\hspace{3.4em}}\hfill 白话：「先翻后转，等于先\underline{\hspace{1.6em}}转再翻。」\hfill

\vspace{0.7em}
\noindent 想一想：三角形的对应等式是 \(fr=r^2f\)，正方形的这条为什么不一样？
\underline{\hspace{5cm}}

\end{document}